\section{Experimental setup}

\subsection{Datasets}

AIDER version 1.0 (CC BY 4.0) supplies five disaster-scene classes and the
primary synthetic resolution audit~\citep{aiderdata}. Its corrected
train/validation/test counts are 4,499/961/969. Satellite Images of Hurricane
Damage (CC BY 4.0) supplies an independent binary corpus through the
fixed Hugging Face revision recorded in the public data-source ledger; corrected counts are
6,997/997/2,000~\citep{hurricanedata}. CRASAR-U-DROIDs (CC BY 4.0) provides
source coordinates, building-damage labels, measured GSD metadata, and aligned
UAS/satellite products at the fixed released revision~\citep{crasar2024data}.

The AIDER frozen split is corrected only by excluding every member of the two
conflicting-label exact-duplicate groups discovered by the pre-specified hash
audit. The Hurricane split uses the fixed seed-0 70/10/20 index assignment and
the same all-members exclusion rule for cross-split hashes.

Intersection-controlled CRASAR classifier training uses guarded crops from Lancaster Canyon
Gate (Hurricane Harvey) and Summerlin San Carlos (Hurricane Ian). The held-out
paired evaluation uses Harlem Heights and McGregor College Parkway South 1
(Hurricane Ian), plus the 13 and 14 October Mexico Beach products (Hurricane
Michael). A prospective third-event sensitivity pairs Steinhatchee River UAS
and post-event crewed-aircraft products from Hurricane Idalia. Evaluation
geographies do not contribute images, labels or outcomes to training or model
selection.

\subsection{Training}

The primary classifier is the ImageNet-1K-initialised
\texttt{microsoft/resnet-18} checkpoint at the revision recorded in the
environment manifest. We train independent seeds
101, 202 and 303 with AdamW (learning rate \(3\times10^{-4}\), weight decay
\(10^{-4}\)), batch size 64, and class-weighted cross-entropy. For class \(c\),
the loss weight is \(N/(C n_c)\), where \(N\) is the training-set size, \(C\)
the number of classes, and \(n_c\) the class count. The checkpoint with the
highest source-validation macro-F1 is retained. AIDER uses six epochs,
Hurricane Damage five, and intersection-controlled CRASAR 12.

Training first resizes to 256 pixels, then applies a random 224-pixel resized
crop with scale range 0.75--1.0 and a horizontal flip with probability 0.5.
Evaluation uses resize-to-256 and a 224-pixel centre crop. Both use ImageNet
channel normalisation. Synthetic resolution degradation uses bicubic
downsampling and upsampling before these transforms. Runs use Python 3.10.20,
PyTorch 2.12.1, torchvision 0.27.1, timm 1.0.28, Transformers 5.12.1 and
Apple M4 Max Metal Performance Shaders.

\subsection{Architecture and operator replication}

After independent review, we prospectively froze a second-architecture and
degradation-operator replication before training or test evaluation. Three
ImageNet-initialised \texttt{timm} MobileNetV3-Small-100 models use the same
clean AIDER split, seeds, optimiser, batch size and six-epoch schedule.
Bicubic, bilinear, nearest-neighbour and box down/up sampling are fixed.

For each architecture, seed and operator, errors are defined at \(s=8\). The
anchor-matched fine-reference score compares \(s=8\) with
\(\{1,2,4\}\), while received-image consistency compares \(s=8\) with
\(\{16,32,64\}\). We report 2,000-replicate paired intervals. The original
torchvision MobileNet weight endpoint returned HTTP 503 before any training;
the pre-evaluation addendum records the switch to
\texttt{timm==1.0.28}.

\subsection{Metrics and fixed analysis}

Error detection is evaluated with AUROC, area under the precision--recall curve
(AUPRC), and false-positive rate at 95\% error recall (FPR95). Classification
metrics include accuracy, balanced accuracy and macro-F1. Threshold-transfer
metrics are coverage, absolute target-coverage error, selective risk and
false-critical rate.

Paired score differences use 10,000 within-seed bootstrap replicates. The
measured-GSD analysis additionally uses the paired spatial-cluster bootstrap
defined above. Seed-level values and sample standard deviations are reported;
model-seed variation is not substituted for building- or site-sampling
uncertainty.

\subsection{Pre-specified criteria}

Criterion A1 requires positive privileged-minus-received mean AUROC on both
synthetic corpora, a positive paired bootstrap interval for every seed, and six
positive seed differences. W11 requires pooled held-out UAS accuracy above the
pooled majority baseline for all seeds and mean balanced accuracy of at least
0.60. W12 requires positive mean received-consistency lift on satellite views,
positive seed-specific spatial-cluster intervals, and positive mean lift at
three or more of four sites. W13 requires the mean satellite absolute coverage
error to be no larger for received consistency than confidence and the mean
received-consistency selective risk to be no higher at target coverage 0.70.
Requiring both inequalities for every seed is reported only as a post-hoc
robustness sensitivity.

The third-event E1 criterion requires mean held-out Idalia UAS balanced
accuracy of at least 0.55. E2 requires complete metrics, selection counts,
physical footprints, raw arrays and overlap-cluster uncertainty for all three
seeds; no ranking-sign or accuracy-direction criterion is imposed.

M1 requires positive bicubic MobileNet anchor-matched gaps and intervals for
all seeds. M2 extends the sign requirement to all four operators, with at least
two positive intervals per operator. M3 requires positive ResNet means and at
least two positive intervals for every corpus/operator combination. M4 requires
aligned fine-reference AUROC to exceed the mean of 100 correspondence-masked
permutations with empirical probability at most 0.05 for all 24
ResNet corpus/seed/operator combinations.
